% ============================================================
%  Mathematiques Appliquees a l'Informatique — Preambule commun
%  Support de cours et Fascicule de travaux pratiques
%  Compilation : tectonic (XeTeX) ou pdflatex, deux passes si refs
% ============================================================
\documentclass[11pt,a4paper]{report}

% ==== Encodage et langue ====
\usepackage{iftex}
\ifPDFTeX
  \usepackage[utf8]{inputenc}
  \usepackage[T1]{fontenc}
  \usepackage{lmodern}
\fi
\usepackage[french]{babel}
\frenchbsetup{StandardLists=true}

% ==== Geometrie ====
\usepackage[a4paper,left=2.4cm,right=2.4cm,top=2.7cm,bottom=2.7cm]{geometry}

% ==== Mathematiques ====
\usepackage{amsmath,amssymb,mathtools}
\usepackage{amsthm} % environnement de demonstration ; boites ci-dessous a compteurs dedies

% ==== Couleurs ====
\usepackage{xcolor}
\definecolor{robleu}{HTML}{16456B}
\definecolor{robleuf}{HTML}{EAF1F7}
\definecolor{roteal}{HTML}{0E7C7B}
\definecolor{rotealf}{HTML}{E9F4F3}
\definecolor{roor}{HTML}{B03A2E}
\definecolor{roorf}{HTML}{FBEDEB}
\definecolor{rogold}{HTML}{B7791F}
\definecolor{rogoldf}{HTML}{FBF3E4}
\definecolor{rovert}{HTML}{2E7D32}
\definecolor{rovertf}{HTML}{EDF6ED}
\definecolor{rogris}{HTML}{4A4A4A}
\definecolor{rogrisf}{HTML}{F3F3F3}

% ==== Graphiques ====
\usepackage{tikz}
\usetikzlibrary{automata,positioning,arrows.meta,calc,shapes.geometric,decorations.pathreplacing,babel,matrix}
\usepackage{pgfplots}
\pgfplotsset{compat=1.15}
\tikzset{badgeRO/.style={rounded corners=1.2mm, draw=#1, line width=0.9pt,
  text=#1, font=\bfseries\small, inner xsep=5mm, inner ysep=1.8mm}}

% ==== Tableaux ====
\usepackage{booktabs,array,multirow,longtable,colortbl}
\newcolumntype{L}[1]{>{\raggedright\arraybackslash}p{#1}}
\newcolumntype{C}[1]{>{\centering\arraybackslash}p{#1}}
\renewcommand{\arraystretch}{1.25}

% ==== Listes ====
\usepackage{enumitem}
\setlist[itemize]{label=\textbullet, itemsep=2pt, topsep=4pt, parsep=2pt}
\setlist[enumerate]{itemsep=2pt, topsep=4pt, parsep=2pt}

% ==== Legendes ====
\usepackage{caption}
\captionsetup{font=small, labelfont=bf}

% ==== Boites pedagogiques ====
\usepackage{tcolorbox}
\tcbuselibrary{skins,breakable}
\usepackage{etoolbox}

\tcbset{robase/.style={breakable, enhanced, boxrule=0.9pt, arc=1.6mm,
  left=3mm, right=3mm, top=2.2mm, bottom=2.2mm, fonttitle=\bfseries,
  before skip=12pt, after skip=14pt}}

\newtcolorbox[auto counter, number within=chapter]{definition}[1][]
  {robase, colback=robleuf, colframe=robleu,
   title={Définition~\thetcbcounter\ifstrempty{#1}{}{~: #1}}}

\newtcolorbox[auto counter, number within=chapter]{propriete}[1][]
  {robase, colback=robleuf!55, colframe=robleu!80!black,
   title={Propriété~\thetcbcounter\ifstrempty{#1}{}{~: #1}}}

\newtcolorbox[auto counter, number within=chapter]{theoreme}[1][]
  {robase, colback=robleu!8, colframe=robleu, boxrule=1.2pt,
   title={Théorème~\thetcbcounter\ifstrempty{#1}{}{~: #1}}}

\newtcolorbox{methode}[1][]
  {robase, colback=rotealf, colframe=roteal,
   title={Méthode\ifstrempty{#1}{}{~: #1}}}

\newtcolorbox[auto counter, number within=chapter]{exemple}[1][]
  {robase, colback=rotealf!70, colframe=roteal,
   title={Exemple résolu~\thetcbcounter\ifstrempty{#1}{}{~: #1}}}

\newtcolorbox{remarque}
  {robase, colback=rogrisf, colframe=rogris!70, boxrule=0.6pt,
   title={Remarque}}

\newtcolorbox{attention}
  {robase, colback=roorf, colframe=roor,
   title={Attention}}

% Boites d'ouverture et de cloture de chapitre (format du modele ESGAE)
\newtcolorbox{quelqueslignes}
  {robase, colback=robleuf!70, colframe=robleu!75!black,
   title={En quelques lignes}}

\newtcolorbox{aretirer}
  {robase, colback=rovertf, colframe=rovert,
   title={Ce qu'il faut retenir}}

\newtcolorbox[auto counter, number within=chapter]{exercice}[1][]
  {robase, colback=rogoldf, colframe=rogold,
   title={Exercice~\thetcbcounter\ifstrempty{#1}{}{~: #1}}}

\newtcolorbox{corrige}[1]
  {robase, colback=rovertf, colframe=rovert,
   title={Corrigé de l'exercice~#1}}

\newcommand{\niveau}[2]{%
\begin{tcolorbox}[enhanced, boxrule=0.9pt, arc=1.6mm, colback=rogoldf,
  colframe=rogold, left=3mm, right=3mm, top=2mm, bottom=2mm,
  title={Public visé et prérequis}, fonttitle=\bfseries,
  before skip=6pt, after skip=14pt]
\textbf{Public concerné :} #1\par\smallskip
\textbf{Prérequis :} #2
\end{tcolorbox}}

% ==== Code informatique ====
\usepackage{listings}
\lstdefinestyle{stylepy}{
  language=Python,
  basicstyle=\ttfamily\footnotesize,
  keywordstyle=\color{robleu}\bfseries,
  commentstyle=\color{rovert}\itshape,
  stringstyle=\color{roor},
  numbers=left, numberstyle=\tiny\color{rogris}, numbersep=8pt,
  frame=single, rulecolor=\color{rogris!50},
  backgroundcolor=\color{rogrisf},
  breaklines=true, breakatwhitespace=false,
  showstringspaces=false, tabsize=4, captionpos=b,
  xleftmargin=16pt, framexleftmargin=14pt,
  morekeywords={True, False, None, self},
  literate={é}{{\'e}}1 {è}{{\`e}}1 {ê}{{\^e}}1 {à}{{\`a}}1
           {ç}{{\c{c}}}1 {î}{{\^i}}1 {ô}{{\^o}}1 {û}{{\^u}}1
           {ù}{{\`u}}1 {É}{{\'E}}1 {È}{{\`E}}1 {À}{{\`A}}1
           {Ô}{{\^O}}1 {Ç}{{\c{C}}}1 {œ}{{\oe}}1 {Œ}{{\OE}}1
           {·}{{\textperiodcentered}}1}
\lstdefinestyle{stylealgo}{
  language={},
  basicstyle=\ttfamily\footnotesize,
  keywordstyle=\color{robleu}\bfseries,
  commentstyle=\color{rovert}\itshape,
  stringstyle=\color{roor},
  numbers=left, numberstyle=\tiny\color{rogris}, numbersep=8pt,
  frame=single, rulecolor=\color{rogris!50},
  backgroundcolor=\color{rogrisf},
  breaklines=true, showstringspaces=false, tabsize=4, captionpos=b,
  xleftmargin=16pt, framexleftmargin=14pt,
  morekeywords={Algorithme, DEBUT, FIN, SI, ALORS, SINON, FINSI, TANTQUE,
    FINTANTQUE, POUR, FINPOUR, FAIRE, FONCTION, FINFONCTION, RETOURNER,
    Entrée, Sortie, Initialiser, Afficher},
  literate={é}{{\'e}}1 {è}{{\`e}}1 {ê}{{\^e}}1 {à}{{\`a}}1
           {ç}{{\c{c}}}1 {î}{{\^i}}1 {ô}{{\^o}}1 {û}{{\^u}}1
           {ù}{{\`u}}1 {É}{{\'E}}1 {È}{{\`E}}1 {À}{{\`A}}1
           {Ô}{{\^O}}1 {Ç}{{\c{C}}}1 {·}{{\textperiodcentered}}1}
\lstdefinestyle{styleout}{
  language={},
  basicstyle=\ttfamily\footnotesize,
  frame=single, rulecolor=\color{rogris!50},
  backgroundcolor=\color{white},
  breaklines=true, showstringspaces=false, xleftmargin=8pt,
  literate={é}{{\'e}}1 {è}{{\`e}}1 {à}{{\`a}}1 {ô}{{\^o}}1 {ç}{{\c{c}}}1 {·}{{\textperiodcentered}}1}

% ==== Titres de chapitres ====
\usepackage{titlesec}
\titleformat{\chapter}[display]
  {\normalfont\bfseries\color{robleu}}
  {\Large\chaptertitlename~\thechapter}{10pt}{\huge}
\titlespacing*{\chapter}{0pt}{-24pt}{26pt}
\titleformat{\section}
  {\normalfont\Large\bfseries\color{robleu}}{\thesection}{0.8em}{}
\titleformat{\subsection}
  {\normalfont\large\bfseries\color{roteal}}{\thesubsection}{0.8em}{}

% ==== En-tetes et pieds de page ====
\usepackage{fancyhdr}
\newcommand{\DocTitre}{Mathématiques Appliquées à l'Informatique}
\pagestyle{fancy}
\fancyhf{}
\fancyhead[L]{\small\itshape\color{rogris} \DocTitre}
\fancyhead[R]{\small\itshape\color{rogris}\nouppercase{\leftmark}}
\fancyfoot[C]{\thepage}
\renewcommand{\headrulewidth}{0.4pt}
\renewcommand{\chaptermark}[1]{\markboth{\chaptername\ \thechapter.~#1}{}}
\fancypagestyle{plain}{\fancyhf{}\fancyfoot[C]{\thepage}%
  \renewcommand{\headrulewidth}{0pt}}

% ==== Divers ====
\usepackage{microtype}
% Interlettrage des pages de garde : propre a pdfTeX, neutralise sans lui
\ifPDFTeX\else
  \ifdefined\textls
    \renewcommand{\textls}[2][]{#2}
  \else
    \newcommand{\textls}[2][]{#2}
  \fi
\fi
\setcounter{secnumdepth}{2}
\setcounter{tocdepth}{1}
\newcommand{\N}{\mathbb{N}}
\newcommand{\Z}{\mathbb{Z}}
\newcommand{\Q}{\mathbb{Q}}
\newcommand{\R}{\mathbb{R}}
\newcommand{\D}{\mathbb{D}}
\newcommand{\bits}{\{0,1\}}
% signe degre, utilisable en math et en texte (\degres est pris par babel-french)
\newcommand{\dgr}{\ensuremath{^{\circ}}}

% ==== Identite du document ====
% (personnalisation : modifier ces macros uniquement)
\newcommand{\Auteur}{Ounimborbitibou DJABON}
\newcommand{\AuteurQualite}{Titulaire de deux Masters : Sciences Mathématiques (AIMS Ghana)\\ et Physique Théorique (Université de Lomé)}
\newcommand{\AuteurLigneB}{Chercheur en Information et Communication Quantique}
\newcommand{\AuteurContact}{+228 92 39 97 21 \ \textbullet \ \ djabon@aims.edu.gh \ \textbullet \ \ github.com/djabelo712}
\newcommand{\Etablissement}{Institut Supérieur CPTEC, Lomé (Togo)}
\newcommand{\InstitutNom}{Institut Supérieur CPTEC}
\newcommand{\InstitutLigne}{Agréé par arrêté N° 00.50/2024 \textbullet{} Avedji, Lomé (Togo) \textbullet{} +228 70 78 78 18 \textbullet{} 96 51 05 50 \textbullet{} cptec228togo@gmail.com}
\newcommand{\AnneeAcad}{Année académique 2026--2027}
\newcommand{\DocKicker}{Support de cours}
\newcommand{\DocSousTitre}{Des fondements logiques au programme opérationnel}
\newcommand{\DocDescription}{Cours complet avec exemples corrigés et code commenté}
\newcommand{\BadgeA}{BTS}
\newcommand{\BadgeB}{Licence}
\newcommand{\BadgeC}{Master Professionnel}
\newcommand{\BadgeUniq}{BTS}
\newcommand{\DocVolumeH}{Volume horaire : 60 heures}

% ==== Page de garde ====
\newcommand{\pagedegarde}{%
\begin{titlepage}
\begin{tikzpicture}[remember picture, overlay]
% ================= bandeau superieur =================
\fill[robleu] (current page.north west)
  rectangle ([yshift=-12.2cm]current page.north east);
\fill[rogold] ([yshift=-12.2cm]current page.north west)
  rectangle ([yshift=-12.42cm]current page.north east);
% filigrane : reseau de noeuds en haut a droite
\begin{scope}
  \draw[white, opacity=0.12, line width=1pt]
    ([xshift=13.5cm,yshift=-2.0cm]current page.north west) --
    ([xshift=15.9cm,yshift=-3.3cm]current page.north west) --
    ([xshift=18.6cm,yshift=-2.2cm]current page.north west) --
    ([xshift=17.7cm,yshift=-5.1cm]current page.north west) --
    ([xshift=14.9cm,yshift=-6.2cm]current page.north west) --
    ([xshift=13.1cm,yshift=-4.5cm]current page.north west) -- cycle;
  \draw[white, opacity=0.12, line width=1pt]
    ([xshift=13.5cm,yshift=-2.0cm]current page.north west) --
    ([xshift=17.7cm,yshift=-5.1cm]current page.north west);
  \draw[white, opacity=0.12, line width=1pt]
    ([xshift=15.9cm,yshift=-3.3cm]current page.north west) --
    ([xshift=14.9cm,yshift=-6.2cm]current page.north west);
  \draw[white, opacity=0.12, line width=1pt]
    ([xshift=18.6cm,yshift=-2.2cm]current page.north west) --
    ([xshift=19.9cm,yshift=-7.2cm]current page.north west);
  \foreach \px/\py in {13.5/-2.0, 15.9/-3.3, 18.6/-2.2, 17.7/-5.1, 14.9/-6.2, 13.1/-4.5, 19.9/-7.2}{%
    \fill[white, opacity=0.18] ([xshift=\px cm,yshift=\py cm]current page.north west)
      circle (0.10);}
\end{scope}
% ligne d'appel
\node[anchor=west, text=rogold!55!white, font=\small]
  at ([xshift=2.1cm,yshift=-2.05cm]current page.north west)
  {\textls[180]{\scshape\DocKicker}};
% titre principal sur deux lignes
\node[anchor=west, text=white, font=\fontsize{38}{45}\selectfont\bfseries]
  at ([xshift=2.0cm,yshift=-4.3cm]current page.north west) {MATHÉMATIQUES};
\node[anchor=west, text=white, font=\fontsize{38}{45}\selectfont\bfseries]
  at ([xshift=2.0cm,yshift=-6.15cm]current page.north west) {GÉNÉRALES};
\fill[rogold] ([xshift=2.05cm,yshift=-7.05cm]current page.north west)
  rectangle ++(2.8,-0.09);
% sous-titre et description
\node[anchor=west, text=white, font=\large\itshape]
  at ([xshift=2.1cm,yshift=-7.8cm]current page.north west) {\DocSousTitre};
\node[anchor=west, text=white!82!robleu, font=\normalsize]
  at ([xshift=2.1cm,yshift=-8.7cm]current page.north west) {\DocDescription};
% pastilles de niveaux
\tikzset{badgeBlanc/.style={rounded corners=1.2mm, draw=white!75!robleu,
  line width=0.8pt, text=white, font=\bfseries\small,
  inner xsep=4.5mm, inner ysep=1.8mm}}
\node[badgeBlanc, anchor=west] (badA)
  at ([xshift=2.1cm,yshift=-10.0cm]current page.north west) {\BadgeUniq};
\node[anchor=west, text=white, font=\small]
  at ([xshift=0.65cm]badA.east) {\DocVolumeH};
% ================= zone centrale : trois panneaux =================
\begin{scope}[shift={(current page.north west)}]
% panneau 1 : equations
\fill[white, draw=rogris!45, line width=0.7pt, rounded corners=2.5mm]
  (1.8,-13.8) rectangle (7.0,-20.05);
\fill[robleu] (2.08,-14.05) rectangle (6.72,-14.11);
\node[text=rogris, font=\footnotesize\scshape] at (4.4,-19.68)
  {\textls[110]{Équations}};
\begin{scope}[shift={(4.4,-16.85)}]
  \node[font=\footnotesize, text=rogris] at (0,0.5) {$x^{2}-5x+6=0$};
  \node[font=\footnotesize, text=rogris] at (0,-0.1) {$\Delta = 1 > 0$};
  \node[font=\footnotesize\bfseries, text=robleu] at (0,-0.7) {$x_{1}=2$ \quad $x_{2}=3$};
\end{scope}
% panneau 2 : systemes lineaires
\fill[white, draw=rogris!45, line width=0.7pt, rounded corners=2.5mm]
  (7.9,-13.8) rectangle (13.1,-20.05);
\fill[roteal] (8.18,-14.05) rectangle (12.82,-14.11);
\node[text=rogris, font=\footnotesize\scshape] at (10.5,-19.68)
  {\textls[110]{Systèmes}};
\begin{scope}[shift={(10.5,-16.85)}]
  \node[font=\footnotesize, text=rogris, align=center] at (0,0.5)
    {$\begin{cases}2x+y=8\\ x-2y=-1\end{cases}$};
  \node[font=\footnotesize, text=rogris] at (0,-0.25) {Cramer : $D=-5$};
  \node[font=\footnotesize\bfseries, text=roteal] at (0,-0.85) {$S=\{(3;\ 2)\}$};
\end{scope}
% panneau 3 : suites numeriques
\fill[white, draw=rogris!45, line width=0.7pt, rounded corners=2.5mm]
  (14.0,-13.8) rectangle (19.2,-20.05);
\fill[rogold] (14.28,-14.05) rectangle (18.92,-14.11);
\node[text=rogris, font=\footnotesize\scshape] at (16.6,-19.68)
  {\textls[110]{Suites}};
\begin{scope}[shift={(16.6,-16.85)}]
  \node[font=\footnotesize, text=rogris] at (0,0.5) {$2;\ 5;\ 8;\ 11;\ \dots\ (+3)$};
  \node[font=\footnotesize, text=rogris] at (0,-0.2) {$3;\ 6;\ 12;\ 24;\ \dots\ (\times 2)$};
  \node[font=\footnotesize\bfseries, text=rogold] at (0,-0.85) {arithmétique et géométrique};
\end{scope}
\end{scope}
% ================= bandeau institut =================
\node[text=robleu, font=\fontsize{13}{15}\selectfont\bfseries]
  at ([yshift=-20.75cm]current page.north) {\InstitutNom};
\node[text=rogris, font=\footnotesize]
  at ([yshift=-21.38cm]current page.north) {\InstitutLigne};
% ================= carte d'auteur =================
\fill[robleuf!60] ([yshift=-21.9cm]current page.north west)
  rectangle (current page.south east);
\draw[robleu, line width=2.2pt]
  ([yshift=-21.9cm]current page.north west) -- ([yshift=-21.9cm]current page.north east);
\draw[rogold, line width=0.9pt]
  ([yshift=-22.14cm]current page.north west) -- ([yshift=-22.14cm]current page.north east);
% nom
\node[text=robleu, font=\fontsize{21}{25}\selectfont\bfseries]
  at ([yshift=-23.4cm]current page.north) {\Auteur};
\fill[rogold] ([xshift=-1.3cm,yshift=-24.32cm]current page.north) rectangle ++(2.6,-0.07);
% lignes de qualite
\node[text width=16.6cm, align=center, text=rogris, font=\small]
  at ([yshift=-25.3cm]current page.north) {\AuteurQualite};
\node[text=rogris, font=\small]
  at ([yshift=-26.4cm]current page.north) {\AuteurLigneB};
\node[text=rogris!80!black, font=\footnotesize]
  at ([yshift=-27.3cm]current page.north) {\AuteurContact};
% ================= bandeau inferieur =================
\fill[robleu] ([yshift=-28.75cm]current page.north west)
  rectangle (current page.south east);
\node[text=white, font=\small]
  at ([yshift=-29.25cm]current page.north) {\textls[150]{\scshape\AnneeAcad}};
\end{tikzpicture}
\end{titlepage}}

% ==== Liens (en dernier) ====
\usepackage[colorlinks=true, linkcolor=robleu, urlcolor=roteal,
            citecolor=robleu]{hyperref}
